% =====================================================================
\chapter{Giới thiệu bài toán}\label{chap:gioithieu}
% =====================================================================

\section{Bối cảnh}
Tại các đô thị lớn như Hà Nội, số lượng ô tô và xe máy tăng nhanh hơn nhiều so với
diện tích bãi đỗ. Người lái thường phải đi vòng nhiều lần để tìm chỗ trống, gây ùn
tắc cục bộ và lãng phí nhiên liệu. Phần lớn bãi xe hiện nay vẫn vận hành thủ công:
vé giấy, ghi biển số bằng tay, tính tiền bằng ước lượng, không biết trước còn bao
nhiêu chỗ và hầu như không có cơ chế cảnh báo an toàn khi xảy ra cháy hoặc ngập
nước — hai rủi ro thường gặp ở tầng hầm.

Một \textbf{bãi đỗ xe thông minh} tự động hóa các công việc này: cảm biến phát hiện
chỗ trống, thẻ RFID kiểm soát ra vào và tính phí, thanh chắn tự mở/đóng, cảm biến
môi trường cảnh báo sự cố và mọi thông tin được hiển thị trực tiếp tại bãi lẫn trên
điện thoại của người quản lý.

\section{Mục tiêu của bài toán}
Đề tài xây dựng một mô hình bãi xe 5 chỗ, 2 cổng, đáp ứng các mục tiêu:
\begin{enumerate}
  \item \textbf{Giám sát chỗ đỗ theo thời gian thực}: biết chính xác chỗ nào trống,
        chỗ nào có xe; màn hình tại bãi cập nhật trong khoảng 0,1 giây, trang web trong
        khoảng 1 giây.
  \item \textbf{Kiểm soát ra vào bằng thẻ RFID}: thẻ quẹt ở cổng vào mở thanh chắn
        vào, cùng thẻ quẹt ở cổng ra mở thanh chắn ra; thanh chắn tự đóng sau 3 giây;
        quẹt sai cổng hoặc bãi đầy bị từ chối và thông báo rõ trên màn hình.
  \item \textbf{Tính phí tự động}: 5.000\,VNĐ cho mỗi giờ bắt đầu, tối thiểu 1 giờ;
        hiển thị hóa đơn khi xe ra.
  \item \textbf{An toàn}: phát hiện lửa — bật còi, hiển thị cảnh báo, \emph{mở cả hai
        thanh chắn} để sơ tán; theo dõi mực nước và nhiệt độ.
  \item \textbf{Giám sát và điều khiển từ xa}: trang web tại \code{http://parking.local}
        dùng được trên mọi điện thoại, không cần cài ứng dụng, có nhật ký xe vào/ra.
  \item \textbf{Hoạt động độc lập}: tự tạo mạng Wi-Fi riêng, không cần Internet hay
        router; tự phục hồi khi mất kết nối.
\end{enumerate}

\section{Ý nghĩa}
\paragraph{Ý nghĩa thực tiễn.} Mô hình chứng minh rằng một bãi xe nhỏ (chung cư mini,
văn phòng, trường học) có thể được tự động hóa bằng linh kiện phổ thông giá rẻ.
Người lái biết ngay còn chỗ hay không; người quản lý thu phí chính xác, có nhật ký
xe ra vào, được cảnh báo cháy/ngập kịp thời; khi cháy, thanh chắn tự mở giúp sơ tán
nhanh.

\paragraph{Ý nghĩa học thuật.} Đây là bài toán nhúng điển hình, vận dụng gần như toàn
bộ kiến thức của môn học: lập trình đa tác vụ với \textbf{hệ điều hành thời gian
thực FreeRTOS}, đồng bộ giữa các tác vụ, điều khiển ngoại vi bằng phần cứng chuyên
dụng (PWM, SPI có DMA, ADC, GPIO), giao tiếp không dây, quản lý bộ nhớ hạn chế
(vài trăm KB RAM), xử lý dữ liệu không tin cậy từ mạng và kiểm thử phần mềm nhúng
khi không có phần cứng.

\section{Phạm vi}
Hệ thống là \textbf{mô hình} quy mô phòng thí nghiệm: 5 chỗ đỗ, 2 cổng, xe mô hình,
thẻ RFID MIFARE. Thanh toán điện tử, nhận dạng biển số và lưu trữ đám mây nằm ngoài
phạm vi và được nêu ở phần hướng phát triển (Chương~\ref{chap:ketluan}).

\section{Kết quả đạt được (tóm tắt)}
\begin{itemize}
  \item Ba bo mạch hoạt động phối hợp, 10 use case (Chương~\ref{chap:usecase}) đều
        được hiện thực và có phép kiểm thử tương ứng.
  \item Firmware máy chủ 1,08\,MB, dùng 36\% RAM tĩnh; còn 29\% phân vùng ứng dụng.
  \item 250 phép kiểm thử tự động đều đạt (Chương~\ref{chap:kiemthu}).
\end{itemize}

\section{Cấu trúc báo cáo}
\begin{table}[H]
\centering\small
\caption{Cấu trúc báo cáo theo yêu cầu của môn học}\label{tab:cautruc}
\begin{tabularx}{\linewidth}{@{}lXl@{}}
\toprule
\textbf{Chương} & \textbf{Nội dung} & \textbf{Yêu cầu} \\ \midrule
\ref{chap:gioithieu} & Giới thiệu bài toán, mục tiêu, ý nghĩa & Giới thiệu (1đ) \\
\ref{chap:usecase} & 10 use case: tác nhân, luồng chính, ngoại lệ, kiểm thử & Use case (10đ) \\
\ref{chap:thanhphan} & Các thành phần phần cứng và phần mềm, thông số quan trọng & Thành phần (2đ) \\
\ref{chap:caidat} & Cài đặt môi trường, hệ điều hành, phần mềm mô phỏng & Cài đặt (1đ) \\
\ref{chap:thietke} & Thiết kế, phân tích kiến trúc, luồng hoạt động, demo & Thiết kế (3đ) \\
\ref{chap:huongdanchay} & Hướng dẫn chạy mã nguồn từng bước & Chạy mã nguồn (2đ) \\
\ref{chap:kiemthu} & Kiểm thử tự động và kịch bản kiểm thử phần cứng & Bổ trợ thiết kế/demo \\
\ref{chap:ketluan} & Tổng kết, hạn chế, hướng phát triển & Kết luận (1đ) \\
\ref{chap:donggop} & Đóng góp của thành viên, phân chia điểm & Đóng góp \\
\bottomrule
\end{tabularx}
\end{table}
